% Primer informe OII464. Se usa el paquete y la portada oficiales de la
% Escuela de Ingeniería Informática PUCV, facilitados para este proyecto.
\documentclass[12pt,letterpaper]{report}
\usepackage[spanish,es-nodecimaldot]{babel}
\usepackage{pucv_inf_2024}
\usepackage{amssymb,booktabs,array,microtype,csquotes}
\usepackage[hidelinks]{hyperref}
\addbibresource{referencias.bib}
\renewcommand{\thesection}{\arabic{section}}
\renewcommand{\thesubsection}{\thesection.\arabic{subsection}}
\renewcommand{\thetable}{\arabic{table}}
\setlength{\headheight}{14pt}

\begin{document}
% Adaptación de la portada para asignaturas de la plantilla PUCV 2024.
\title{Programación de cirugías electivas\\[0cm]
mediante búsqueda tabú y selección\\[0.2cm]
adaptativa de movimientos}
\author{José Basualto, Cristopher Dauros y Carlos Ruz}
\date{2 de octubre de 2026}

\begin{titlepage}
  \makeatletter
  \thispagestyle{empty}
  \begin{picture}(0,0)
    \put(-75,-80){\includegraphics[width=20cm]{Portadas/imagenes/encabezado.png}}
  \end{picture}
  \begin{flushright}
    \vspace*{\fill}
    {\textcolor{titlecolor}{\LARGE\bfseries\MakeUppercase{\\[3cm]\@title}}}
    \vspace*{\fill}

    {\textcolor{titlecolor}{\Large\bfseries\scshape
      \begin{tabular}{@{}r@{}}
        José Basualto\\[0.2cm]
        Cristopher Dauros\\[0.2cm]
        Carlos Ruz
      \end{tabular}}}\\[1cm]
    {\textcolor{titlecolor}{\bfseries\textsc{
      Primer informe\\[0.2cm]
      OII464: Desarrollo de Diseños Híbridos para Optimización\\[0.2cm]
      \@date}}}
  \end{flushright}
  \makeatother
\end{titlepage}


\pagenumbering{roman}
\chapter*{Resumen}
La programación de cirugías electivas exige decidir el orden de los pacientes
y coordinar quirófanos, anestesiólogos y cirujanos en dos etapas. Las salas
pueden diferir entre anestesia y cirugía. La espera entre etapas debe
respetar un máximo estricto. Una mala secuencia prolonga la jornada y retiene
recursos.
El desafío es reducir el tiempo de término de la última cirugía y esa espera
sin relajar las restricciones clínicas y operativas.

Proponemos una búsqueda tabú para explorar intercambios, inserciones y
cambios de sala. Un bandido multibrazo dinámico (D-MAB) aprende durante la
corrida qué movimiento conviene proponer a partir de las mejoras observadas;
un detector de cambios reinicia sus estimaciones si el rendimiento cae. Las
Un planificador estricto repara y evalúa las soluciones.
La comparación incluye Tabu con selección uniforme, una ablación sin detector
y tres reglas constructivas, con igual presupuesto de evaluaciones.

Una validación preliminar con veinte semillas nuevas en doce instancias
sintéticas produjo mejoras y empeoramientos de magnitud semejante: no
demuestra una ventaja general del controlador. El bloqueo mediano fue cero,
lo que limita la capacidad del catálogo para distinguir métodos en esa
métrica. El aporte es un protocolo reproducible para estudiar cuándo
conviene adaptar movimientos.

\noindent\textbf{Palabras clave:} programación de cirugías; búsqueda tabú;
selección adaptativa de operadores; aprendizaje en línea; bloqueo.

\clearpage
\pagenumbering{arabic}
\section{Introducción}
En la programación intrajornada de cirugías electivas se asignan recursos
compartidos a operaciones con precedencia. El caso docente de OII464 contiene
una etapa de anestesia y otra de cirugía por paciente. Cada una tiene
duración, preparación, transición, limpieza y conjuntos de salas y personal
elegible. Si el quirófano de cirugía no está disponible al terminar la
anestesia, la primera sala y su profesional pueden permanecer retenidos hasta
el traslado. Una solución debe evitar superposiciones de recursos y respetar
la espera máxima de cada paciente. Minimizar la suma de duraciones no
representaría el problema: las duraciones son datos, mientras que el orden y
las asignaciones determinan cuándo termina la jornada. Por eso examinamos el
\emph{makespan} y, por separado, la espera entre etapas como indicador de
bloqueo.

La presentación del taller propone integrar una metaheurística y mejorar
reglas clásicas, además de reducir el makespan y evitar quirófanos bloqueados
\parencite{lara2026}. Para estudiar una arquitectura híbrida hace falta una
comparación controlada: cambiar a la vez la regla de factibilidad, la
representación o el presupuesto impediría atribuir diferencias al componente
de aprendizaje. También importa que los movimientos discretos no tengan la
misma utilidad durante toda la búsqueda. Una mezcla fija de intercambio,
inserción y cambios de sala no responde a esa variación; un controlador en
línea podría hacerlo, aunque su exploración y sus reinicios pueden asimismo
perjudicar una corrida. La pregunta es si D-MAB mejora \emph{de manera
consistente} la calidad de Tabu Search frente a seleccionar los mismos
movimientos uniformemente, manteniendo constantes el decodificador, el
presupuesto y las soluciones iniciales.

Los trabajos relacionados motivan esta comparación, pero no la resuelven.
\textcite{akbarzadeh2024} incorporan varias fases de recursos y predicción de
factibilidad en secuenciación quirúrgica; \textcite{yin2026} usan búsqueda
tabú con varios operadores entre hospitales. La elección aprendida de
operadores también tiene precedentes: \textcite{yu2024} emplean
Q-learning en programación y reprogramación de cirugías, mientras
\textcite{niroumandrad2024} guían Tabu con clasificadores en otra tarea
clínica. \textcite{dacosta2008} proponen D-MAB para selección adaptativa,
combinando UCB1 y detección Page--Hinkley, y \textcite{dolatkhah2026} ya
aplican D-MAB a operadores genéticos en planificación de quirófanos. Esos
aportes impiden reivindicar como novedad general el uso de aprendizaje o de
D-MAB en este dominio. La brecha concreta de nuestro estudio es evaluar su
integración \emph{dentro de Tabu}, con espera máxima estricta, salas de ambas
etapas independientes y una comparación emparejada sobre el caso del curso.

El objetivo de esta primera entrega es especificar la representación, los
componentes metaheurístico y de aprendizaje, su interacción y un protocolo
reproducible para medir calidad, tiempo y robustez. El repositorio docente
\parencite{saicooh2026} aporta contexto y datos sintéticos; se trabajó en un
repositorio independiente. La presentación menciona cuatro archivos JSON,
pero los experimentos disponibles aquí usan doce instancias YAML de la copia
local del catálogo. Esta distinción y la función objetivo exacta se declaran
para evitar comparaciones impropias con puntajes históricos.

\section{Propuesta metodológica}
\subsection{Problema, decisiones y decodificación}
Sea $J$ el conjunto de pacientes. Una solución propuesta
$S=(\pi,\rho^{A},\rho^{C})$ incluye una permutación $\pi$ y una sala
elegible para anestesia y otra, potencialmente distinta, para cirugía de
cada paciente. El planificador asigna personal elegible y ubica cada etapa
en el primer instante compatible con la disponibilidad de sala y profesional.
Para el paciente $j$, sean $F^A_j$ el fin de anestesia incluida su limpieza,
$T^C_j$ el inicio efectivo de cirugía y $F^C_j$ el fin de cirugía incluida
su limpieza. Definimos
\begin{equation}
  W_j=\max\{0,T^C_j-F^A_j\},\qquad
  0\leq W_j\leq L_j,\qquad
  C_{\max}=\max_{j\in J}F^C_j,
  \label{eq:restricciones}
\end{equation}
donde $L_j$ es la espera máxima de la instancia. El límite se verifica como
\emph{restricción dura}: una propuesta que lo viola no puede ser la mejor
solución. Además se exige elegibilidad y ausencia de superposición de salas
y profesionales. Cuando las etapas usan salas distintas, la sala y el
profesional de anestesia permanecen ocupados hasta $T^C_j$; si comparten
sala, ésta sigue ocupada durante la cirugía y limpieza. $W_j$ mide la espera
entre etapas usada como \emph{proxy} de bloqueo, no toda la indisponibilidad
de un quirófano.

La codificación de claves aleatorias genera preferencias de orden y de dos
salas por paciente. Un decodificador repara esas preferencias: si el paciente
preferido no es programable en ese momento, prueba los siguientes; para el
elegido busca un par de salas factible que se aparte lo menos posible de las
preferencias. El personal se asigna de manera voraz entre los elegibles.
Por ello el movimiento solicitado puede convertirse en otra solución tras
la reparación. El espacio realmente explorado depende de este decodificador
y no comprende necesariamente todos los calendarios factibles. A diferencia
del piloto inicial, no se fuerza equilibrio entre pabellones ni la misma sala
para ambas etapas.

Minimizamos el criterio fijado antes de la validación,
\begin{equation}
  f(S)=C_{\max}+10^{-6}\sum_o t_o^{\mathrm{ini}}
       +0{,}5\sum_{j\in J}W_j+1{,}4\max_{j\in J}W_j,
  \label{eq:objetivo}
\end{equation}
donde $t_o^{\mathrm{ini}}$ es el inicio de reserva del recurso para
cada operación, incluida la preparación. El término $10^{-6}$ desempata
soluciones próximas. Reportamos $C_{\max}$ y las esperas por separado: una
mejora de $f$ no implica que todas sus componentes disminuyan. La
presentación docente considera también equilibrio de carga
\parencite{lara2026}; lo excluimos porque una asignación desigual puede
minimizar mejor el tiempo y la espera. Por tanto, $f$ no equivale al puntaje
histórico del taller.

\subsection{Componente metaheurístico: búsqueda tabú}
La búsqueda comienza desde la mejor de treinta soluciones evaluadas con el
mismo planificador: FIFO, SPT, LPT y veintisiete generadas por claves
aleatorias. Las reglas clásicas forman parte de la inicialización, de modo
que superarlas al final está garantizado por construcción y no demuestra por
sí solo el aporte de Tabu o D-MAB. En cada ronda se generan quince candidatos
distintos mediante cinco tipos de movimiento: intercambiar dos pacientes,
extraer e insertar uno en otra posición, cambiar su sala de anestesia,
cambiar la de cirugía o cambiar ambas. Con probabilidad $0{,}65$ se elige un
paciente entre los tres con mayor espera actual; en otro caso se elige del
conjunto completo. El sesgo concentra parte de la búsqueda donde hay
congestión, sin impedir exploración global.

La memoria tabú conserva las últimas siete soluciones para desincentivar
retornos. Una aspiración permite aceptar una solución tabú si mejora el
mejor valor global. El estado local puede empeorar para salir de una zona de
búsqueda, pero el mejor calendario observado se guarda por separado. Las
tres políticas comparadas comparten exactamente estos movimientos, la
reparación, las soluciones iniciales y un máximo de $3.030$ evaluaciones
estrictas, incluidas las treinta iniciales.

\subsection{Componente de aprendizaje y flujo}
El aprendizaje se realiza \emph{online} en cada corrida, sin entrenamiento
previo ni un modelo contextual. Antes de generar cada candidato, D-MAB
escoge un movimiento $a$ mediante el índice UCB1
\begin{equation}
  a^*=\arg\max_a\left(\bar r_a+
      \sqrt{\frac{2\ln(\sum_b n_b)}{n_a}}\right),
  \label{eq:ucb}
\end{equation}
tras probar inicialmente cada brazo. Aquí $n_a$ es su número de observaciones
y $\bar r_a$ su recompensa media. Después de evaluar el candidato se entrega
una recompensa en $[0,1]$:
\begin{equation}
  r_a=\min\!\left\{1,
    \frac{\max\{0,f(S_{\mathrm{local}})-f(S_{\mathrm{cand}})\}}
         {\max\{1,0{,}02 f(S_{\mathrm{inicial}})\}}\right\}.
  \label{eq:recompensa}
\end{equation}
La referencia local se fija al comienzo del lote de quince candidatos. Una
propuesta duplicada recompensa con cero al bandido, aunque no consume una
evaluación estricta. Page--Hinkley detecta caídas de recompensa por brazo
con $\delta=0{,}01$ y umbral $\lambda=0{,}35$; cuando dispara, reinicia
las estimaciones de \emph{todos} los movimientos, pero conserva la solución,
la memoria tabú y el mejor valor. D-MAB aprende qué \emph{propuesta más
reparación} fue útil; no diagnostica directamente un óptimo local.

El flujo \emph{offline} consiste en verificar instancias, definir función
objetivo, parámetros, baselines y semillas. En línea se repite: seleccionar
movimiento, generar vecino, reparar, comprobar factibilidad, evaluar,
actualizar bandido y estado tabú. Como ablación, UCB1 usa la misma selección
sin Page--Hinkley; Tabu mixto selecciona uno de los cinco movimientos con
probabilidad uniforme. Así la comparación principal estima el efecto de la
selección aprendida y la ablación aísla el reinicio dinámico.

\subsection{Datos, evaluación y alcance de la evidencia}
El catálogo contiene doce instancias sintéticas en YAML: tres réplicas
de cada tamaño, con 15, 20, 25 y 30 cirugías. Se verifican
esquema y hashes, y se registran tamaño, salas, personal, duraciones y límites
de espera. Son características descriptivas para análisis posterior, no
entradas del bandido. Al crecer el número de cirugías también crecen los
recursos, por lo que una curva agrupada por tamaño no aísla el efecto del
número de pacientes. Ninguna instancia contiene datos clínicos reales.

El piloto usó semillas $0$--$19$ para fijar la arquitectura. La validación
emparejada usó semillas $20$--$39$ en \emph{las mismas doce instancias}: tres
políticas por instancia y semilla, $720$ corridas y $2.181.600$
evaluaciones estrictas presupuestadas. Comparamos diferencias por pareja
instancia--semilla, además de medianas de objetivo, makespan, espera total y
máxima, porcentaje factible, tiempo y curvas de convergencia ilustrativas.
La mediana del tiempo por corrida fue cercana a $0{,}9$ segundos. Una
reejecución de 36 corridas reconstruyó calendarios factibles y reprodujo los
valores guardados. La comparación frente a FIFO, SPT y LPT se muestra como
referencia constructiva, con la salvedad de que estas reglas ya siembran la
búsqueda.

\begin{table}[ht]
  \centering
  \caption{Comparación preliminar frente a Tabu mixto en 240 pares
  instancia--semilla. Un cambio de hasta 0,25 min se considera similar.}
  \label{tab:resultados}
  \begin{tabular}{lrrr}
    \toprule
    Política & Mejora & Similar & Empeora\\
    \midrule
    Tabu + UCB1 & 105 & 34 & 101\\
    Tabu + D-MAB & 106 & 39 & 95\\
    \bottomrule
  \end{tabular}
\end{table}

El efecto de D-MAB cambia de signo entre instancias y la mediana de sus doce
diferencias medianas respecto de Tabu mixto es aproximadamente cero. La
ablación directa tampoco muestra una ventaja consistente del reinicio
Page--Hinkley. El bloqueo total mediano fue cero en todas las instancias y
políticas; este catálogo discrimina poco en esa métrica y las diferencias
provienen principalmente del makespan. Las semillas nuevas permiten estudiar
la variación aleatoria del algoritmo \emph{dentro del catálogo conocido},
pero no prueban generalización a instancias no vistas ni optimalidad global.
Como siguiente prueba se requieren instancias sintéticas de estrés con más
congestión y un conjunto de instancias reservado antes de ajustar parámetros.

\noindent\textbf{Amenazas de validez y trabajo siguiente.} El presupuesto
igualado cuenta llamadas al planificador estricto, no propuestas repetidas ni
segundos de procesador. En efecto, una propuesta duplicada puede actualizar
el bandido con recompensa cero sin gastar una evaluación, y la política
uniforme no registra ese intento con la misma granularidad. La reparación
puede modificar el movimiento escogido; por tanto, una preferencia aprendida
por un brazo refleja la utilidad de proponerlo \emph{junto con} su reparación,
no el efecto puro de una operación de vecindario. Una evaluación posterior
deberá registrar tasas de reparación, elecciones y recompensas por operador,
y comparar además presupuestos de tiempo. También se requieren intervalos de
incertidumbre agrupados por instancia: las veinte semillas de una misma
réplica comparten los datos y no constituyen veinte problemas independientes.
Finalmente, nuevas instancias con esperas más exigentes permitirán observar
si el componente aprendido aporta valor precisamente donde hay riesgo de
bloqueo.

\section{Conclusión}
El problema exige coordinar secuencia, salas y personal sin violar la espera
máxima entre anestesia y cirugía. La propuesta combina Tabu Search con D-MAB
para escoger movimientos según la recompensa observada durante cada corrida.
La validación inicial no demuestra superioridad general del controlador y
señala que el catálogo actual apenas discrimina bloqueo. El valor de la
arquitectura reside en que deja aislado y reproducible el efecto del
aprendizaje; las siguientes entregas deberán probarlo en instancias no
observadas y con mayor congestión.

\clearpage
\printbibliography[title={Referencias}]
\end{document}
